\documentstyle[12pt,fullpage,steve]{article}
\setcounter{page}{1}

\begin{document}

\title{EE 150 Midterm \#2}
\author{Nov 13, 1996}
\date{}
\maketitle

{\bf Read First:}
\begin{itemize}
\item {\bf **READ** THE DIRECTIONS CAREFULLY!}

\item This exam is open book, open notes, and closed neighbor!

\item You have a maximum 50 minutes to complete the exam.  Hand in your paper when you are done.
We will pick up all midterms and leave the room at 55 minutes after the hour.  

\item There are 3 problems worth a total of 100 points plus a Bonus question worth 5 points.  
Problem I is worth 3 * 15 or 45 marks, Problem II is worth 40 marks, and
Problem III is worth 15 marks.   Be careful not to spend too much time on
any particular problem.  {\bf If you get stuck, go on to the next
problem.}  You can always come back to the one that is giving you
problems.  

\item As a guideline, allow about 15 minutes for each of problems I II and III,
which leaves you a couple of minutes at the end to check your work.
To ensure you have read these instructions, you get two more bonus points
simply for writing your account name on the bottom of the last page of this
exam.
\end{itemize}

% ---------------------------------------------------------------------------
\pagebreak
{\bf PROBLEM \#1} (45 points, 15 points each)
 
Once again, show (below each program) what each of the following program
fragments do.  To determine the answer, work through the code by hand and show
us the output that results.  {\em You {\bf must}  show how these programs
execute by using the  ``table'' model we have seen repeatedly in class in order
to get full credit. }  

Hint:  The table method of showing the output involves showing the time or
execution steps on one axis (say the horizontal axis) and each variable that
changes on the other axis (say the vertical axis).     You should show the
output somewhere in the same table.

\begin{itemize}
\item {\bf Part 1-A}

\begin{verbatim}
main()
{
  int t, v = 0;
  for (v = 10; t < 6, t > 5; v++)
    {
      printf(" In Loop: t= %i v= %i\n", t, v);  
      t += 2;
      v -= 3;
    }
  printf("Done: v = %i, t = %i\n", v, t);
  return 0;
}
\end{verbatim}

% ---------------------------------------------------------------------------
\pagebreak

\item {\bf Part 1-B}

\begin{verbatim}
main()
{
  int f1 (int);
  int j = 12;
  int y = 15;
  do
    {
      int y = 10;
      y = f1 ( y );
      printf("In Loop: y = %i, j = %i \n", y, j), --j;
    }
  while( y > 0 && j != 10 );
  return 0;
}

int f1 ( int y ){  return y -= y; }
\end{verbatim}

%---------------------------------------------------------------------------
\pagebreak
\item {\bf Part 1-C}

\begin{verbatim}
main()
{
  int i, a = 0;
  int A[5] = { 4, 3, 2, 1, 0};
  for (i = 0; i < 5; i = (a += 1))
    switch( i )
      {
      case 1: A[i]= 5;  break;
      case 2: 
      case 3: A[i] = A[ (int) i/2 ];   break;
      case 4: A[i] = 7; break;
      default: A[i] = A[i + 3]; break;
      }
  for (i = 0; i < 5; i++)
    printf("Array A element %i is %i\n", i, A[i]); 

  return 0;
}
\end{verbatim}

\end{itemize}

%---------------------------------------------------------------------------
\pagebreak

{\bf PROBLEM \#2} (40 points)

Given a function, {\tt writeT}, {\em that has already
been written for you } .  {\tt writeT} takes two arguments,
a character, c, and a number, n, and writes out character c exactly (int)(n/3)
times. For example, {\tt writeT('x', 9)} writes {\tt xxx}.
<P>
Using this function, {\em write a new function } , {\tt writeManyT},
where the call
\begin{verbatim}
int stuff[5] = {9, 12, 8, 16, 1};
writeT(stuff, 5, 6, '=');
\end{verbatim}

produces the following output:  
\begin{verbatim}
===xxx===0
===xxxx===0
===xx===2
===xxxxx===1
\end{verbatim}

Note that the final integer on each of the preceeding lines can be calculated
as  {\bf stuff[i] - ( (int)( stuff[i] / 3) ) * 3}.  {\bf You may assume that
all integer values in stuff are positive.} 

A sketch of the function and its expected parameters is shown below.
Your job is to fill in its body.

\begin{verbatim}
/*
   Display table's parameters:

     a -    An array of integers.  For each array element, x, we want to 
            printout (int)(x/3) characters 'x'.
     n  -   The number of entries in the array
     m  -   We may assume that m is even.  We want to printout (m/2) instances 
            of character c before and after the  x's. 
     c-     The character to write out as per m above
*/

writeManyT(int a[], int n, int m, char c)
{
... write your code here!
}
\end{verbatim}

% ---------------------------------------------------------------------------
\pagebreak

{\bf Write Your Solution to Problem \#2 on this Page}


% ---------------------------------------------------------------------------
\newpage

{\bf PROBLEM \#3} (15 points)

For each of the following program fragments, show an alternative (i.e. expanded
or more straight-forward) form:
\begin{verbatim}
1. switch(i)
        {
            case 1: printf("ok");      break;
            case 2: 
            case 3: printf("no way");  break;
           default: printf("ho vay");  break;
        }






\end{verbatim}

\begin{verbatim}
2. i++;   /* do not use post or pre-fix operators */

\end{verbatim}

\begin{verbatim}
3. A ? x++ : x--;


\end{verbatim}

\begin{verbatim}
4. int a, x = 0;

\end{verbatim}

\begin{verbatim}
5. if (a *= b, a *= c, a > 10)   /* do not use comma operators */
      printf("fred");
   else
      printf("sam");     




\end{verbatim}

% ---------------------------------------------------------------------------
\pagebreak

{\bf Bonus PROBLEM \#4} (+5 extra points)

Write a short paragraph about how your group will be/is  attacking the slide-5
group problem, and in particular what your contribution is going to be.

\end{document}